\section{Pierwsze twierdzenie o niezupełności}\label{sec:pierwsze-twierdzenie}
\begin{definicja}[$\omega$-niesprzeczność]\label{res:9:1}
Teoria $T$ jest \emph{$\omega$-niesprzeczna}, jeśli nie istnieje formuła $\varphi(x)$ taka, że $T\vdash\exists x\,\varphi(x)$ oraz jednocześnie $T\vdash\neg\varphi(\overline n)$ dla każdego $n\in\mathbb N$. Każda teoria $\omega$-niesprzeczna jest niesprzeczna. Istotnie, w teorii sprzecznej można udowodnić zarówno $\exists x\,(x=x)$, jak i każde $\neg(\overline n=\overline n)$.
\end{definicja}

\begin{twierdzenie}[Pierwsze twierdzenie Gödla, wersja Gödla]\label{res:9:2}
Niech $T\supseteq Q$ będzie efektywnie aksjomatyzowalna, ze standardowym predykatem z Lematu~\ref{res:7:7}. Zbudujmy $G$ jako punkt stały formuły $\neg\operatorname{Bew}_T(x)$, czyli
\begin{equation}\label{eq:9:1}
Q\vdash G\leftrightarrow\neg\operatorname{Bew}_T(\ulcorner G\urcorner).
\end{equation}
Wówczas:
\begin{enumerate}[label=(\alph*)]
\item jeśli $T$ jest niesprzeczna, to $T\nvdash G$;
\item jeśli $T$ jest $\omega$-niesprzeczna, to $T\nvdash\neg G$.
\end{enumerate}
W szczególności, jeśli $T$ jest $\omega$-niesprzeczna, to $G$ jest zdaniem nierozstrzygalnym w $T$, więc $T$ jest niezupełna.
\end{twierdzenie}
\begin{proof}
Istnienie $G$ gwarantuje Lemat~\ref{res:8:1} zastosowany do $\psi(x):=\neg\operatorname{Bew}_T(x)$. Intuicyjnie $G$ głosi: „nie jestem dowodliwe w $T$”.

\emph{(a) $T\nvdash G$.} Przypuśćmy przeciwnie, że $T\vdash G$, i wybierzmy pełny kod $p$ tego dowodu. Relacja $\operatorname{Prov}_T(p,\ulcorner G\urcorner)$ zachodzi w $\mathbb N$, więc Lemat~\ref{res:7:7} daje $Q\vdash\operatorname{Prf}_T(\overline p,\ulcorner G\urcorner)$, skąd $T\vdash\operatorname{Bew}_T(\ulcorner G\urcorner)$. Z~\eqref{eq:9:1} i $T\vdash G$ wynika natomiast $T\vdash\neg\operatorname{Bew}_T(\ulcorner G\urcorner)$. Teoria dowodziłaby zdania i jego negacji, wbrew niesprzeczności. Zatem $T\nvdash G$.

\emph{(b) $T\nvdash\neg G$.} Teraz zakładamy $\omega$-niesprzeczność, a więc także niesprzeczność. Punkt (a) mówi, że żadne standardowe $n$ nie jest kodem dowodu $G$. Lemat~\ref{res:7:7} daje zatem osobno dla każdego $n$
\begin{equation}\label{eq:9:2}
T\vdash\neg\operatorname{Prf}_T(\overline n,\ulcorner G\urcorner).
\end{equation}
Gdyby $T\vdash\neg G$, z~\eqref{eq:9:1}, przez klasyczną eliminację podwójnej negacji, otrzymalibyśmy
\begin{equation}\label{eq:9:3}
T\vdash\exists p\,\operatorname{Prf}_T(p,\ulcorner G\urcorner).
\end{equation}
Zdanie~\eqref{eq:9:3} wraz z całą rodziną~\eqref{eq:9:2} jest dokładnie sytuacją zabronioną przez Definicję~\ref{res:9:1}, dla formuły $\varphi(p)=\operatorname{Prf}_T(p,\ulcorner G\urcorner)$. Zatem $T\nvdash\neg G$.

Nie zastąpiliśmy rodziny~\eqref{eq:9:2} jednym wewnętrznym twierdzeniem z kwantyfikatorem uniwersalnym. Właśnie dlatego ten argument potrzebuje $\omega$-niesprzeczności, a nie tylko zwykłej niesprzeczności.
\end{proof}

\begin{uwaga}[Prawdziwość zdania Gödla]\label{res:9:3}
Jeśli $T$ jest niesprzeczna, to $G$ jest prawdziwe w $\mathbb N$: z punktu (a) nie istnieje rzeczywisty dowód $G$, więc $\mathcal N\models\neg\operatorname{Bew}_T(\ulcorner G\urcorner)$. Równoważność~\eqref{eq:9:1} jest dowodliwa już w $Q$, a $\mathcal N\models Q$, toteż zachodzi także w $\mathbb N$. Nie zakładamy tutaj prawdziwości wszystkich aksjomatów $T$. Do punktu (b) wystarczyłaby nawet $\Sigma_1$-poprawność $T$ (dowodzenie tylko prawdziwych zdań $\Sigma_1$), ponieważ~\eqref{eq:9:3} jest fałszywym zdaniem tej klasy. Sama niesprzeczność nie uprawnia jednak do punktu (b) dla tego samego $G$.
\end{uwaga}

\begin{twierdzenie}[Wersja Rossera --- zwykła niesprzeczność]\label{res:9:4}
Każda niesprzeczna, efektywnie aksjomatyzowalna teoria $T\supseteq Q$ jest niezupełna: istnieje zdanie $R$ takie, że $T\nvdash R$ i $T\nvdash\neg R$.
\end{twierdzenie}
\begin{proof}
\emph{Poprawna formuła Rossera.} Funkcja $\operatorname{neg}(f)$ zwracająca kod negacji jest p.r.; dla niepoprawnych kodów ma wartość zastępczą $0$. Nie jest ona symbolem języka arytmetyki, dlatego wybieramy formułę $N(f,z)$ reprezentującą ją w $Q$ (Tw.~\ref{res:7:6}) i definiujemy
\[
\operatorname{Bew}_T^{R}(f):=\exists z\,\exists p\,
\bigl(N(f,z)\land\operatorname{Prf}_T(p,f)\land\forall q<p\,\neg\operatorname{Prf}_T(q,z)\bigr).
\]
Lemat diagonalny daje zdanie $R$ z $Q\vdash R\leftrightarrow\neg\operatorname{Bew}_T^R(\ulcorner R\urcorner)$. Niech $r=\ulcorner R\urcorner$, $r^-=\ulcorner\neg R\urcorner$. Ponieważ $Q\vdash\forall z\,(N(\overline r,z)\leftrightarrow z=\overline{r^-})$, w $Q$ otrzymujemy konkretną równoważność
\begin{equation}\label{eq:9:4}
R\leftrightarrow\neg\exists p\,
\bigl(\operatorname{Prf}_T(p,\overline r)\land
\forall q<p\,\neg\operatorname{Prf}_T(q,\overline{r^-})\bigr).
\end{equation}
\emph{Pierwszy przypadek:} załóżmy $T\vdash R$. Niech $a$ będzie pełnym kodem dowodu $R$. Z niesprzeczności wynika $T\nvdash\neg R$, więc żadne $q<a$ nie jest kodem dowodu $\neg R$. Lemat~\ref{res:7:7} i skończony odcinek z Lematu~\ref{res:4:4} dają
\[
Q\vdash\operatorname{Prf}_T(\overline a,\overline r),\qquad
Q\vdash\forall q<\overline a\,\neg\operatorname{Prf}_T(q,\overline{r^-}).
\]
Wprowadzając świadka $p=\overline a$, dostajemy w $T$ zaprzeczenie prawej strony~\eqref{eq:9:4}, czyli $T\vdash\neg R$: sprzeczność. Zatem $T\nvdash R$.

\emph{Drugi przypadek:} załóżmy $T\vdash\neg R$. Niech $b$ będzie pełnym kodem tego dowodu. Wiemy już, że dowodu $R$ nie ma; potrzebujemy jednak tylko skończenie wielu przypadków $p\le b$. Mamy
\[
Q\vdash\operatorname{Prf}_T(\overline b,\overline{r^-}),\qquad
Q\vdash\forall p\le\overline b\,\neg\operatorname{Prf}_T(p,\overline r).
\]
Pracujemy teraz w $Q$ z dowolnym $p$. Z~\eqref{eq:4:2} wynika $p\le\overline b\lor\overline b<p$. W pierwszej gałęzi odpada $\operatorname{Prf}_T(p,\overline r)$. W drugiej $q=\overline b$ obala $\forall q<p\,\neg\operatorname{Prf}_T(q,\overline{r^-})$. Wobec tego $Q$ dowodzi negacji koniunkcji pod kwantyfikatorem w~\eqref{eq:9:4} dla każdego $p$, a więc dowodzi $R$. To przeczy $T\vdash\neg R$ i niesprzeczności $T$.

Nie przenieśliśmy nieskończonej rodziny twierdzeń o numerałach pod jeden kwantyfikator. Jeden konkretny dowód $\neg R$ wyznaczył skończoną granicę $b$ i obsłużył wszystkich pozostałych kandydatów naraz.
\end{proof}

\begin{wniosek}[Prawdziwa arytmetyka nie jest rekurencyjnie przeliczalna]\label{res:9:5}
Zbiór $\operatorname{Th}(\mathcal N)$ wszystkich zdań prawdziwych w $\mathcal N$ nie jest rekurencyjnie przeliczalny.
\end{wniosek}
\begin{proof}
Gdyby był rekurencyjnie przeliczalny, można by wziąć go za efektywny zbiór aksjomatów teorii zawierającej $Q$. Teoria ta byłaby niesprzeczna, ponieważ $\mathcal N$ byłoby jej modelem, i zupełna, ponieważ każde zdanie lub jego negacja jest prawdziwe w $\mathcal N$. Przeczy to Twierdzeniu~\ref{res:9:4}.
\end{proof}

\begin{wniosek}[Nierozstrzygalność dowodliwości]\label{res:9:6}
Jeśli $T\supseteq Q$ jest niesprzeczna i efektywnie aksjomatyzowalna, to zbiór kodów zdań dowodliwych w $T$ nie jest rekurencyjny.
\end{wniosek}
\begin{proof}
Niech przeciwnie rozstrzygalna relacja $D(n)$ oznacza „$n$ jest kodem zdania dowodliwego w $T$”. Wybieramy jej dwustronną reprezentację $\rho_D(x)$ w $Q$ i punkt stały $Q\vdash A\leftrightarrow\neg\rho_D(\ulcorner A\urcorner)$. Jeśli $T\vdash A$, reprezentacja daje $T\vdash\rho_D(\ulcorner A\urcorner)$, wbrew równoważności i niesprzeczności. Jeśli $T\nvdash A$, reprezentacja daje $Q\vdash\neg\rho_D(\ulcorner A\urcorner)$, a stąd $T\vdash A$: znów sprzeczność.
\end{proof}

